\documentclass[11pt]{report}
\input{preambolo-verifica.tex}

\titleformat{\subsection}[block]
  {\normalfont\large\bfseries\color{cyan}}{\thesubsection}{1em}{}
\titlespacing*{\subsection}{0pt}{0.6\baselineskip}{0.2\baselineskip}

\newcommand{\eserc}[1]{\bigskip\noindent\textbf{Esercizio #1.}\quad}

\begin{document}
\chapter*{Correzione — Verifica sulle equazioni di II grado incomplete}

\section*{Fila A}

\subsection*{Esercizi (fino a 9 punti)}

\eserc{1} $\dfrac{1-\sqrt2}{\sqrt[4]2}\,n^2=0$

Il coefficiente $\dfrac{1-\sqrt2}{\sqrt[4]2}$ è diverso da zero (perché $1-\sqrt2\neq0$), quindi deve essere $n^2=0$, da cui $n=0$.
\[
\boxed{n=0}
\]

\eserc{2} $-\sqrt\pi\,x^2-\dfrac{15}{4}=0$

\begin{align*}
-\sqrt\pi\,x^2 &= \dfrac{15}{4}\\
x^2 &= -\dfrac{15}{4\sqrt\pi}
\end{align*}
Il secondo membro è negativo, ma $x^2$ non può mai essere negativo: l'equazione è impossibile.
\[
\boxed{\text{nessuna soluzione reale}}
\]

\eserc{3} $9-2z^2=-9$

\begin{align*}
-2z^2 &= -9-9=-18\\
z^2 &= 9
\end{align*}
\[
\boxed{z=\pm3}
\]

\eserc{4} $(2x+3)^2=5x+9$

\begin{align*}
4x^2+12x+9 &= 5x+9\\
4x^2+12x+9-5x-9 &=0\\
4x^2+7x &=0\\
x(4x+7)&=0
\end{align*}
\[
\boxed{x=0 \ \text{oppure}\ x=-\dfrac74}
\]

\eserc{5} $\left(\dfrac x2-2\right)\left(\dfrac x2+2\right)=8$

Differenza di quadrati:
\begin{align*}
\dfrac{x^2}{4}-4 &= 8\\
\dfrac{x^2}{4} &= 12\\
x^2 &= 48
\end{align*}
\[
\boxed{x=\pm\sqrt{48}=\pm4\sqrt3}
\]

\eserc{6} $-\dfrac{y+3}{4} - (3y-2)^2 = 2y - \dfrac{19}{4}$

Moltiplico tutto per $4$:
\begin{align*}
-(y+3) - 4(3y-2)^2 &= 8y-19\\
-y-3-4(9y^2-12y+4) &= 8y-19\\
-y-3-36y^2+48y-16 &= 8y-19\\
-36y^2+47y-19 &= 8y-19\\
-36y^2+47y-8y-19+19 &=0\\
-36y^2+39y &=0\\
-3y(12y-13)&=0
\end{align*}
\[
\boxed{y=0 \ \text{oppure}\ y=\dfrac{13}{12}}
\]

\subsection*{Esercizi bonus}

\eserc{7} $-3(x-1)^5=(x-1)^3$

Porto tutto a sinistra e raccolgo $(x-1)^3$:
\begin{align*}
-3(x-1)^5-(x-1)^3 &=0\\
(x-1)^3\left[-3(x-1)^2-1\right] &=0
\end{align*}
Quindi $(x-1)^3=0 \Rightarrow x=1$, oppure $-3(x-1)^2-1=0 \Rightarrow (x-1)^2=-\dfrac13$, impossibile.
\[
\boxed{x=1}
\]

\eserc{8} $x^2+6x+9=7$

Il primo membro è già il quadrato di un binomio:
\begin{align*}
(x+3)^2 &= 7\\
x+3 &= \pm\sqrt7
\end{align*}
\[
\boxed{x=-3\pm\sqrt7}
\]

\eserc{9} $x\left(\dfrac {x^2}2 +1\right)(x^3-2)^4\left(2x+\dfrac14\right)=0$

Un prodotto è nullo se e solo se lo è (almeno) uno dei fattori:
\begin{itemize}
\item $x=0$
\item $\dfrac{x^2}2+1=0 \Rightarrow x^2=-2$, impossibile
\item $(x^3-2)^4=0 \Rightarrow x^3=2 \Rightarrow x=\sqrt[3]2$
\item $2x+\dfrac14=0 \Rightarrow x=-\dfrac18$
\end{itemize}
\[
\boxed{x=0,\quad x=\sqrt[3]2,\quad x=-\dfrac18}
\]

\clearpage
\section*{Fila B}

\subsection*{Esercizi (fino a 9 punti)}

\eserc{1} $-\sqrt3\,t^2-\dfrac{22}{9}=0$

\begin{align*}
-\sqrt3\,t^2 &= \dfrac{22}{9}\\
t^2 &= -\dfrac{22}{9\sqrt3}
\end{align*}
Il secondo membro è negativo: l'equazione è impossibile.
\[
\boxed{\text{nessuna soluzione reale}}
\]

\eserc{2} $13-4x^2=-3$

\begin{align*}
-4x^2 &= -3-13=-16\\
x^2 &= 4
\end{align*}
\[
\boxed{x=\pm2}
\]

\eserc{3} $\dfrac{\sqrt3-2}{\sqrt[3]5}\,m^2=0$

Il coefficiente $\dfrac{\sqrt3-2}{\sqrt[3]5}$ è diverso da zero (perché $\sqrt3\neq2$), quindi $m^2=0$.
\[
\boxed{m=0}
\]

\eserc{4} $(3t+2)^2=7t+4$

\begin{align*}
9t^2+12t+4 &= 7t+4\\
9t^2+12t+4-7t-4 &=0\\
9t^2+5t &=0\\
t(9t+5)&=0
\end{align*}
\[
\boxed{t=0 \ \text{oppure}\ t=-\dfrac59}
\]

\eserc{5} $(3t-2)(3t+2)=32$

\begin{align*}
9t^2-4 &= 32\\
9t^2 &= 36\\
t^2 &= 4
\end{align*}
\[
\boxed{t=\pm2}
\]

\eserc{6} $-\dfrac{x+4}{5} - (2x-3)^2 = 3x - \dfrac{49}{5}$

Moltiplico tutto per $5$:
\begin{align*}
-(x+4) - 5(2x-3)^2 &= 15x-49\\
-x-4-5(4x^2-12x+9) &= 15x-49\\
-x-4-20x^2+60x-45 &= 15x-49\\
-20x^2+59x-49 &= 15x-49\\
-20x^2+59x-15x-49+49 &=0\\
-20x^2+44x &=0\\
-4x(5x-11)&=0
\end{align*}
\[
\boxed{x=0 \ \text{oppure}\ x=\dfrac{11}{5}}
\]

\subsection*{Esercizi bonus}

\eserc{7} $-2(t-1)^7=3(t-1)^5$

Porto tutto a sinistra e raccolgo $(t-1)^5$:
\begin{align*}
-2(t-1)^7-3(t-1)^5 &=0\\
(t-1)^5\left[-2(t-1)^2-3\right] &=0
\end{align*}
Quindi $(t-1)^5=0 \Rightarrow t=1$, oppure $-2(t-1)^2-3=0 \Rightarrow (t-1)^2=-\dfrac32$, impossibile.
\[
\boxed{t=1}
\]

\eserc{8} $t^2+8t+16=9$

\begin{align*}
(t+4)^2 &= 9\\
t+4 &= \pm3
\end{align*}
\[
\boxed{t=-1 \ \text{oppure}\ t=-7}
\]

\eserc{9} $\sqrt[3] y\left(2y+\dfrac14\right)(y^3-2)^3\left(\dfrac {y^2}2 -1\right)=0$

Un prodotto è nullo se e solo se lo è (almeno) uno dei fattori:
\begin{itemize}
\item $\sqrt[3]{y}=0 \Rightarrow y=0$
\item $2y+\dfrac14=0 \Rightarrow y=-\dfrac18$
\item $(y^3-2)^3=0 \Rightarrow y^3=2 \Rightarrow y=\sqrt[3]2$
\item $\dfrac{y^2}2-1=0 \Rightarrow y^2=2 \Rightarrow y=\pm\sqrt2$
\end{itemize}
\[
\boxed{y=0,\quad y=-\dfrac18,\quad y=\sqrt[3]2,\quad y=\pm\sqrt2}
\]

\clearpage
\section*{Fila C}

\subsection*{Esercizi (fino a 9 punti)}

\eserc{1} $7+2x^2=7$

\[
2x^2=0
\]
\[
\boxed{x=0}
\]

\eserc{2} $3y^2-1=11$

\begin{align*}
3y^2 &= 12\\
y^2 &= 4
\end{align*}
\[
\boxed{y=\pm2}
\]

\eserc{3} $x^2=x$

\begin{align*}
x^2-x &=0\\
x(x-1)&=0
\end{align*}
\[
\boxed{x=0 \ \text{oppure}\ x=1}
\]

\eserc{4} $4z^2+5=1$

\[
4z^2=1-5=-4 \quad\Rightarrow\quad z^2=-1
\]
Il quadrato di un numero non può essere negativo: l'equazione è impossibile.
\[
\boxed{\text{nessuna soluzione reale}}
\]

\eserc{5} $(x+3)^2=9+7x$

\begin{align*}
x^2+6x+9 &= 9+7x\\
x^2+6x+9-9-7x &=0\\
x^2-x &=0\\
x(x-1)&=0
\end{align*}
\[
\boxed{x=0 \ \text{oppure}\ x=1}
\]

\eserc{6} $-(x-1)^2=1$

\[
(x-1)^2=-1
\]
Il quadrato di un numero non può essere negativo: l'equazione è impossibile.
\[
\boxed{\text{nessuna soluzione reale}}
\]

\subsection*{Esercizio bonus}

\eserc{7} $-2(x^2+2)(x+2)(x^4+1)(x^2-1)=0$

Il fattore $-2$ non è mai nullo; un prodotto è nullo se e solo se lo è (almeno) uno degli altri fattori:
\begin{itemize}
\item $x^2+2=0 \Rightarrow x^2=-2$, impossibile
\item $x+2=0 \Rightarrow x=-2$
\item $x^4+1=0 \Rightarrow x^4=-1$, impossibile (una potenza pari non può essere negativa)
\item $x^2-1=0 \Rightarrow x=\pm1$
\end{itemize}
\[
\boxed{x=-2,\quad x=1,\quad x=-1}
\]

\end{document}
